\documentclass[10pt,a4paper]{amsart}
\usepackage[margin=19mm]{geometry}
\usepackage{amsmath,amssymb,mathtools}
\usepackage[scr=rsfs,cal=euler]{mathalfa}
\usepackage{xfrac}
\usepackage[unicode,hidelinks,bookmarks=false]{hyperref}
\newcommand{\ie}{i.e.\ }
\newcommand{\cf}{cf.\ }
\newcommand{\resp}{respectively\ }
\newcommand{\Subsection}{\subsection}
\providecommand{\nobreakdash}{\nobreak\hbox{-}}
\setlength{\emergencystretch}{2em}
\multlinegap0pt
\usepackage{bm}
\usepackage{bbm}
\usepackage{dsfont}
\usepackage{xspace}
\usepackage{cancel}
\usepackage{mathtools}
\usepackage{amsthm}
\makeatletter
\@ifundefined{theorem}{\newtheorem{theorem}{Theorem}}{}
\@ifundefined{proposition}{\newtheorem{proposition}[theorem]{Proposition}}{}
\makeatother
\newcommand{\bb}{\mathbb}
\newcommand{\vphi}{\varphi}
\title[$S$-Integral Points in Orbits on $\mathbb{P}^1$]{$S$-Integral Points in Orbits on $\mathbb{P}^1$}
\author{Jit Wu Yap}
\date{Source: arXiv:2312.05094v2 (CC BY 4.0)}
\begin{document}
\maketitle

\noindent\emph{Source and licence.} $S$-Integral Points in Orbits on $\mathbb{P}^1$. Version arXiv:2312.05094v2 (CC BY 4.0), distributed under the Creative Commons Attribution 4.0 International licence, as recorded in the arXiv licence metadata for this exact version and shown on the abstract page https://arxiv.org/abs/2312.05094v2. Full rights notice: licenses/arxiv-2312.05094-CC-BY-4.0.txt.\par\smallskip

\noindent\textit{Editorial scope.} Counted unit: the Proposition whose source label is \texttt{DerivativeBound2} in the article (source0.tex lines 762--770, source label \texttt{DerivativeBound2}), delivered together with the article's own complete original proof of it (source0.tex lines 772--808, one \texttt{proof} environment of 37 lines). No mathematics is written, repaired or re-derived: statement and proof are the article's own text line for line. Required context retained uncounted: DerivativeBound1 (lines 732--737). Every definition, display and label of those lines is retained unchanged. Omitted from this package and not redistributed: 1--731, 738--761, 771--771, 809--1082. Nothing inside a retained excerpt is omitted or altered: every retained body is delivered exactly as the article's lines are written, so no omission receipt of any kind accompanies this package. Citations: the retained excerpts carry no citation command at all, so no bibliography entry of the article is transcribed and no reference is redistributed. Reference adaptation: none was needed and none was made; every reference inside the delivered unit resolves inside it, so no unresolved reference or citation is produced. Printed slips kept literal: no printed word, symbol, hypothesis or number of the counted statement or of its proof was altered. Layout and class: the article's own class is \texttt{amsart} with packages including amsmath, amsthm, hyperref, ragged2e, enumitem, tikz-cd, mathtools, accents, amssymb; the wrapper is the lane's pinned \texttt{amsart} editorial wrapper at 10pt on A4, which restates the theorem-like environments the retained text needs and adds the article's own macro definitions that the retained text needs (\texttt{\textbackslash bb}, \texttt{\textbackslash vphi}), with extra wrapper packages (bm, bbm, dsfont, xspace, cancel, mathtools). The delivered numbering is the unit's own. Assets: no retained span contains a figure environment, a graphics-inclusion command, a PDF-page inclusion, a table or a TikZ picture; no image file is copied into this package, the package declares no asset dependency and no image file is redistributed with it. Licence: version arXiv:2312.05094v2 of this article is distributed under the Creative Commons Attribution 4.0 International licence, as recorded in the arXiv licence metadata for this exact version and shown on the abstract page https://arxiv.org/abs/2312.05094v2. A full rights notice is delivered at licenses/arxiv-2312.05094-CC-BY-4.0.txt. Characters: every retained excerpt is pure ASCII, so the delivered engine, XeLaTeX with its default Latin Modern text font, reports no missing character.
\begin{proposition} \label{DerivativeBound1}
 Assume that $\beta,\vphi(\beta) \in K$. Write 
 $$\vphi(z) = \vphi(\beta) + d_1 (z-\beta) + d_2 (z-\beta)^2 + \cdots$$
 Then there exists $C = C(d) > 0$ depending only on $d$ such that  
$$h(d_n) \leq Cn(h(\vphi) + h(\beta)+1).$$. 
\end{proposition}

\begin{proposition} \label{DerivativeBound2}
Let $K$ be a number field and $\vphi$ a rational map defined over $K$ of degree $d \geq 2$. Assume that $\beta$ and $\vphi(\beta)$ are both in $K$. Then for any place $v \in M_K$, there exists constants $C_1 = C_1 (d) > 0, C_2 = C_2(d) > 0$, both depending only on the degree $d$, with $C_1 \geq C_2$, such that if $z \in \bb{C}_v$ satisfies 
\begin{equation} \label{eq: SuperAttracting2}
\log |z-\beta|_v \leq -C_1[K:\bb{Q}](h(\vphi) + h(\beta)+1),
\end{equation}
then
$$\left| \log |\vphi(z)-\vphi(\beta)|_v - k \log |z-\beta|_v \right| \leq C_2 [K:\bb{Q}](h(\vphi) + h(\beta) + 1)$$
where $k$ is the order of vanishing of $\beta$ for the equation $\vphi(z) = \vphi(\beta)$. 
\end{proposition}

\begin{proof}
For simplicity, let $M = [K;\bb{Q}](h(\vphi) + h(\beta)+1)$. We expand 
$$\vphi(x) = \vphi(\beta) + d_1(z-\beta) + d_2(z-\beta)^2 + \cdots.$$
Let $d_k$ be the smallest $k$ for which $d_k \not = 0$. Then $k \leq 2d-2$ and 
$$h(d_k) \leq C(h(\vphi) + h(\beta)+1)$$
where $C$ is the constant from Proposition \ref{DerivativeBound1}. Thus 
\begin{equation} \label{eq: SuperAttracting6}
\log |d_k|_v \geq - CkM.
\end{equation}
Similarly for all $n \geq k$, we have 
\begin{equation} \label{eq: SuperAttracting5}
\log |d_n|_v \leq Cn M.
\end{equation}
We may now bound 
$$\left|\left(\vphi(z) - \vphi(\beta) \right)  - d_k (z-\beta)^k \right| \leq \sum_{i=k+1}^{\infty} \left|d_i (z-\beta)^i \right|.$$
Now let
$$\log |z-\beta|_v = -C_1 M$$
then
$$|d_i (z-\beta)^i| \leq e^{-i (C_1 - C) M}.$$
Then summing up from $i=k+1$ to $\infty$ gives an upper bound 
$$\sum_{i=k+1}^{\infty} \left|d_i (z-\beta)^i \right| \leq 2 e^{-i(C_1 - C)M}.$$
Now let's assume that $C_1 \geq 3d C$, so that $i(C_1 - C) \geq (k+ \frac{1}{d})C_1$ as $i \geq k+1$ and $k \leq 2d$. Then we obtain 
\begin{equation} \label{eq: SuperAttracting4}
\left|\vphi(z) - \vphi(\beta)| - d_k |z-\beta|^k \right| \leq 2e^{-(k + \frac{1}{d})C_1 M}.
\end{equation}
By \eqref{eq: SuperAttracting5} and \eqref{eq: SuperAttracting6}, we have 
$$e^{(-kC_1 - kC)M} \leq d_k|z-\beta|^k \leq e^{(-kC_1 + Ck)M}.$$
We now assume that $C_1$ is large enough so that
\begin{equation} \label{eq: SuperAttracting7}
e^{-(k+\frac{1}{d})C_1 M} \leq \frac{1}{2} e^{M(-kC_1 - kC)} \leq \frac{1}{2} d_k |z-\beta|^k.
\end{equation}
It follows from \eqref{eq: SuperAttracting4}
$$ \frac{1}{2} |z-\beta|^k \leq |\vphi(z) - \vphi(\beta)| \leq \frac{3}{2} d_k |z-\beta|^k$$
and so
$$\left|\log |\vphi(z) - \vphi(\beta)|_v ^{-1} - k \log |z-\beta|^{-1}_v \right| \leq \log |d_k|_v + \log 2$$
which by \eqref{eq: SuperAttracting5} gives us the $C_2$ we want as desired, where we take $C_1$ to be a constant larger than $3dC$ so that \eqref{eq: SuperAttracting7} holds. 
\end{proof}

\end{document}
